\documentclass[11pt]{article}
\usepackage[utf8]{inputenc}
\usepackage[T1]{fontenc}
\usepackage[italian]{babel}
\usepackage{geometry}
\geometry{a4paper, margin=1in}
\usepackage{xcolor}
\usepackage{listings}
\usepackage{hyperref}

\usepackage{tcolorbox}
\tcbuselibrary{breakable}
\begin{document}

\begin{tcolorbox}[colback=blue!5!white,colframe=blue!75!black,title=\textbf{DOMANDA}, breakable]
Esercizio 1\\
Modificare a piacimento il contenuto del file \texttt{public/index.html} e valutarne l'impatto grafico.
\end{tcolorbox}

\begin{tcolorbox}[colback=green!5!white,colframe=green!75!black,title=\textbf{RISPOSTA}, breakable]
\textbf{Soluzione Esercizio 1:} Modifica grafica e impatto visivo

\textbf{Operazioni Preliminari:}
\begin{enumerate}
    \item Ho copiato tutto il codice sorgente originale in \texttt{es34 - WebSocket Esercizio 1}.
    \item Ho stravolto il file \texttt{public/index.html} e riscritto da zero il foglio di stile \texttt{public/styles.css}.
    \item Ho avviato il server Node.js in background. Per visualizzare il risultato, basta collegarsi a \texttt{http://localhost:4000}.
\end{enumerate}

\textbf{Analisi dell'impatto grafico (Il restyling):}

\textbf{L'interfaccia originale era molto "grezza":} un contenitore bianco su sfondo bianco, pulsanti squadrati e tipografia di sistema, caratteristiche tipiche di un prototipo base. 

Per modernizzare l'interfaccia e offrire un impatto visivo ("WOW factor") all'altezza delle applicazioni web contemporanee, ho implementato le seguenti migliorie strutturali e stilistiche in CSS puro:

\begin{enumerate}
    \item \textbf{Il Tema (Dark Gradient):}\\
    Ho sostituito lo sfondo bianco con un gradiente diagonale scuro (da blu notte profondo a indaco). Questo fa risaltare il riquadro centrale della chat e riduce l'affaticamento visivo (Dark Mode nativa).

    \item \textbf{Glassmorphism (Effetto Vetro Smerigliato):}\\
    Invece di un semplice colore a tinta unita per il container della chat, ho utilizzato uno sfondo semi-trasparente (\texttt{rgba(255, 255, 255, 0.05)}) combinato con la direttiva \texttt{backdrop-filter: blur(12px)}. Questo simula l'effetto di un pannello di vetro sospeso sopra lo sfondo, conferendo una profondità tridimensionale tipica delle UI di iOS o di Windows 11.

    \item \textbf{Tipografia Moderna:}\\
    Ho rimosso il font di default (Nunito/System) e importato \textit{Poppins} da Google Fonts. Le sue geometrie rotonde e i pesi marcati (\texttt{font-weight: 600}) donano alla WebApp un tono molto più amichevole, netto e leggibile.

    \item \textbf{Animazioni e Dettagli (Micro-Interactions):}
    \begin{itemize}
        \item Ho trasformato il bottone "Send" da un blocco rettangolare a un pulsante circolare flottante inserendo un'icona SVG vettoriale minimalista (un aeroplanino di carta).
        \item Ho aggiunto degli effetti di \textit{hover} (il bottone si ingrandisce al passaggio del mouse e si rimpicciolisce al click, animato fluidamente tramite \texttt{transition: all 0.3s ease}).
        \item I nuovi messaggi che appaiono in chat, invece di spuntare di scatto, hanno ora un'animazione \texttt{fadeIn} CSS (\texttt{@keyframes}) che li fa "scivolare" dolcemente dal basso verso l'alto rendendo l'esperienza visiva fluida e gratificante.
    \end{itemize}
\end{enumerate}

Il risultato è una chat che, pur mantenendo l'identica infrastruttura e tecnologia WebSocket dell'esercizio precedente (Node.js e Socket.io), appare all'utente finale come un'applicazione premium, dimostrando quanto l'impatto grafico (Frontend) possa alterare drasticamente la percezione di qualità del software, a parità di backend.
\end{tcolorbox}

\end{document}
